% LQTA-D Article 3 -- Editorial Pass 02
% Cycles, memory geometry, and fixed-carrier transitions candidate v0.3
% This file is written as a replacement candidate for Sections sec:memory and sec:fixed.

\section{Cycles, holonomy, and the geometry of collective temporal memory}
\label{sec:memory}

For any closed oriented cycle $C$, define the CAL period
\begin{equation}
\Pi_C(\gamma)=\sum_{e\in C}\epsilon_{C,e}\gamma_e,
\label{eq:cycle-period}
\end{equation}
where $\epsilon_{C,e}=\pm1$ records agreement or disagreement with the chosen cycle orientation. Under a local recalibration
\begin{equation}
\gamma\longmapsto\gamma+D_0q,
\end{equation}
the period is invariant,
\begin{equation}
\Pi_C(\gamma+D_0q)=\Pi_C(\gamma),
\end{equation}
because the coboundary contribution telescopes around every closed cycle. Thus CAL cycle periods are properties of the relational comparison state modulo local recalibration.

With the frozen unit edge inner product, write
\begin{equation}
\gamma=P_{\mathrm{ex}}\gamma+P_\perp\gamma
      =D_0\phi+\eta,
\qquad
\eta=P_\perp\gamma\in\ker D_0^{\mathsf T}.
\label{eq:hodge-memory-decomposition}
\end{equation}
The vector $\eta$ is gauge invariant under addition of exact cochains and contains exactly the part of CAL that is not representable by a global family of local rates.

\begin{proposition}[Dimension of the memory sector]
\label{prop:memory-dimension}
Let $G$ be a connected finite CAL carrier. Then
\begin{equation}
\dim\ker D_0^{\mathsf T}=\beta_1(G)=|E|-|V|+1.
\label{eq:memory-dimension}
\end{equation}
Consequently, the non-exact CAL sector has one real degree of freedom per independent cycle.
\end{proposition}

\paragraph*{Proof.}
For a connected graph, $\operatorname{rank}D_0=|V|-1$. Rank--nullity for $D_0^{\mathsf T}:\mathbb R^{|E|}\to\mathbb R^{|V|}$ gives
\begin{equation}
\dim\ker D_0^{\mathsf T}=|E|-(|V|-1)=|E|-|V|+1=\beta_1(G).
\end{equation}
\hfill$\square$

\begin{corollary}[Tree rigidity]
\label{cor:tree-rigidity}
If $G$ is a connected tree, then $\beta_1(G)=0$ and therefore
\begin{equation}
\eta=0
\end{equation}
for every CAL state. Every CAL assignment on a tree is integrable.
\end{corollary}

This gives the first structural separation:
\begin{equation}
\boxed{
\text{nonhomogeneity does not require cycles, whereas irreducible CAL memory does.}
}
\label{eq:nonhom-vs-memory}
\end{equation}

Define the local allocation and total magnitude of the non-exact component by
\begin{equation}
\rho_i=\frac12\sum_{e\ni i}\eta_e^2,
\qquad
\Rcal=\sum_i\rho_i=\sum_e\eta_e^2=\lVert\eta\rVert^2.
\label{eq:rho-R}
\end{equation}
The quantity $\Rcal$ is invariant under local recalibration. It measures the magnitude of the non-exact CAL component relative to the frozen cochain inner product. It is not identified with energy, entropy, elapsed time, matter density, action, or spacetime curvature.

It is useful to distinguish the scalar magnitude $\Rcal$ from the full memory configuration $\eta$. States with the same value of $\Rcal$ need not carry the same cycle content when the cycle space has dimension greater than one.

\subsection{One independent cycle: scalar memory up to sign}

Suppose $\beta_1(G)=1$. Choose a unit vector $h$ spanning $\ker D_0^{\mathsf T}$. Then every non-exact state has the form
\begin{equation}
\eta=\alpha h,
\qquad
\Rcal=\alpha^2.
\label{eq:beta1-one-coordinate}
\end{equation}
Hence a nonzero value of $\Rcal$ fixes the non-exact state up to the two possibilities
\begin{equation}
\eta=+\sqrt{\Rcal}\,h
\qquad\text{or}\qquad
\eta=-\sqrt{\Rcal}\,h.
\label{eq:beta1-one-two-points}
\end{equation}
The scalar $\Rcal$ therefore loses only a sign/orientation of cycle content in the one-cycle sector.

\begin{proposition}[Rigidity of source-free endpoints for $\beta_1=1$]
\label{prop:beta1-sourcefree-rigidity}
Let $\eta$ and $\eta'$ be two non-exact states on a fixed connected carrier with $\beta_1=1$. If
\begin{equation}
\lVert\eta'\rVert=\lVert\eta\rVert,
\end{equation}
then
\begin{equation}
\eta'=\eta
\qquad\text{or}\qquad
\eta'=-\eta.
\label{eq:beta1-rigidity}
\end{equation}
If $\eta=0$, the only equal-$\Rcal$ endpoint is $\eta'=0$.
\end{proposition}

\subsection{Two or more independent cycles: a continuous memory manifold}

For $\beta_1(G)=b\ge2$, choose an orthonormal basis $\{h_1,\ldots,h_b\}$ of the cycle subspace. Writing
\begin{equation}
\eta=\sum_{a=1}^{b}x_a h_a,
\end{equation}
gives
\begin{equation}
\Rcal=\sum_{a=1}^{b}x_a^2.
\end{equation}
For every fixed $\Rcal=R>0$, the set of non-exact states is therefore the sphere
\begin{equation}
\mathcal M_R
=
\{\eta\in\ker D_0^{\mathsf T}:\lVert\eta\rVert^2=R\}
\cong S^{b-1}.
\label{eq:memory-sphere}
\end{equation}
Thus the total amount $\Rcal$ does not determine how memory is distributed among independent cycle directions.

Equivalently, if $O$ is any orthogonal transformation of the cycle subspace,
\begin{equation}
\eta' = O\eta
\qquad\Longrightarrow\qquad
\Rcal'=\Rcal.
\label{eq:orthogonal-memory}
\end{equation}
For $b\ge2$ there is a continuous family of such states.

The distinction can be made concrete on a carrier containing two independent edge-disjoint cycles. Let $h_1$ and $h_2$ be their normalized cycle vectors. Then
\begin{equation}
\eta(\theta)=\sqrt{R}\bigl(\cos\theta\,h_1+\sin\theta\,h_2\bigr)
\label{eq:two-cycle-rotation}
\end{equation}
has $\Rcal=R$ for all $\theta$, while its edgewise and hence local $\rho_i$ allocation varies with $\theta$. The total nonintegrable content can therefore be conserved while its relational distribution changes.

This is a purely kinematic statement about the state space. No preferred value of $\theta$, no rotation law, and no autonomous motion on $\mathcal M_R$ is selected by the present ontology.

\section{Fixed-carrier transition classes}
\label{sec:fixed}

Let a realized CAL update occur on a fixed carrier. Write
\begin{equation}
\gamma' = \gamma+\Delta\gamma,
\qquad
\eta' = \eta+\Delta\eta,
\end{equation}
where
\begin{equation}
\Delta\eta=P_\perp\Delta\gamma.
\label{eq:delta-eta}
\end{equation}
Because the projector is fixed, the change in total non-exact content is exactly
\begin{align}
\Delta\Rcal
&=\lVert\eta'\rVert^2-\lVert\eta\rVert^2 \\
&=2\langle\eta,\Delta\eta\rangle+\lVert\Delta\eta\rVert^2.
\label{eq:fixed-transition-law}
\end{align}

\begin{definition}[Fixed-carrier transition classes]
A realized update is called
\begin{enumerate}
\item \emph{exact-only} if $\Delta\eta=0$;
\item \emph{transverse source-free} if $\Delta\eta\neq0$ but $\Delta\Rcal=0$;
\item \emph{source-positive} if $\Delta\Rcal>0$;
\item \emph{source-negative} if $\Delta\Rcal<0$.
\end{enumerate}
\end{definition}

The terminology ``source'' refers only to the structural balance of $\Rcal$. It does not assign an energetic or thermodynamic meaning to the sign.

\begin{proposition}[Integrable-origin positivity]
\label{prop:integrable-origin-positive}
If the initial state is integrable, $\eta=0$, then
\begin{equation}
\Delta\Rcal=\lVert\Delta\eta\rVert^2\ge0.
\label{eq:integrable-origin-law}
\end{equation}
Therefore an integrable state can leave the integrable sector only through a source-positive realized update. A source-free update starting from $\eta=0$ remains integrable.
\end{proposition}

\begin{proposition}[Erasure of nonintegrability]
\label{prop:erasure}
If a realized update sends $\eta$ to $\eta'=0$, then
\begin{equation}
\Delta\Rcal=-\lVert\eta\rVert^2=-\Rcal.
\label{eq:erasure-law}
\end{equation}
Thus complete erasure of pre-existing nonintegrability is necessarily source-negative on fixed support.
\end{proposition}

\subsection{Source-free transitions as motion on an equal-$\Rcal$ sphere}

A transverse update is source-free exactly when
\begin{equation}
\lVert\eta'\rVert=\lVert\eta\rVert.
\label{eq:sourcefree-sphere}
\end{equation}
The geometry therefore depends sharply on the cycle rank.

For $\beta_1=1$, Proposition~\ref{prop:beta1-sourcefree-rigidity} implies that the only endpoints at fixed nonzero $\Rcal$ are
\begin{equation}
\eta'=\eta
\qquad\text{and}\qquad
\eta'=-\eta.
\end{equation}
The second is a nontrivial source-free inversion.

For $\beta_1\ge2$, every orthogonal transformation of the cycle subspace supplies a source-free endpoint,
\begin{equation}
\eta'=O\eta,
\qquad O\in O(\beta_1),
\label{eq:sourcefree-orthogonal}
\end{equation}
and the equal-$\Rcal$ state space contains a continuum. Consequently, conservation of $\Rcal$ is much weaker than conservation of the detailed temporal-memory configuration.

\subsection{Radial transition family}

A particularly transparent family is
\begin{equation}
\eta'=\lambda\eta,
\qquad\lambda\in\mathbb R.
\label{eq:radial-map}
\end{equation}
For $\eta\neq0$,
\begin{equation}
\Delta\Rcal=(\lambda^2-1)\Rcal.
\label{eq:radial-source}
\end{equation}
Hence
\begin{align}
|\lambda|&<1 &&\Longrightarrow &&\Delta\Rcal<0,\\
|\lambda|&=1 &&\Longrightarrow &&\Delta\Rcal=0,\\
|\lambda|&>1 &&\Longrightarrow &&\Delta\Rcal>0.
\end{align}
The cases $\lambda=1$ and $\lambda=-1$ already give two inequivalent source-free maps whenever $\eta\neq0$.

These results are classifications of possible realized changes, not laws selecting one change. In particular, neither the sign of $\Delta\Rcal$ nor conservation of $\Rcal$ provides an endogenous transition principle. The fixed-carrier geometry therefore anticipates the generator no-go of Section~\ref{sec:stop}: structural invariants constrain evolution but do not determine it.
